
\documentclass{article}
\usepackage{colortbl}
\usepackage{makecell}
\usepackage{multirow}
\usepackage{supertabular}

\begin{document}

\newcounter{utterance}

\centering \large Interaction Transcript for game `hot\_air\_balloon', experiment `air\_balloon\_survival\_en\_negotiation\_easy', episode 4 with glm{-}4.6{-}t1.0.
\vspace{24pt}

{ \footnotesize  \setcounter{utterance}{1}
\setlength{\tabcolsep}{0pt}
\begin{supertabular}{c@{$\;$}|p{.15\linewidth}@{}p{.15\linewidth}p{.15\linewidth}p{.15\linewidth}p{.15\linewidth}p{.15\linewidth}}
    \# & $\;$A & \multicolumn{4}{c}{Game Master} & $\;\:$B\\
    \hline

    \theutterance \stepcounter{utterance}  
    & & \multicolumn{4}{p{0.6\linewidth}}{
        \cellcolor[rgb]{0.9,0.9,0.9}{
            \makecell[{{p{\linewidth}}}]{
                \texttt{\tiny{[P1$\langle$GM]}}
                \texttt{You are participating in a collaborative negotiation game.} \\
\\ 
\texttt{Together with another participant, you must agree on a single set of items that will be kept. Each of you has your own view of how much each item matters to you (importance). You do not know how the other participant values the items. Additionally, you are given the effort each item demands.} \\
\texttt{You may only agree on a set if the total effort of the selected items does not exceed a shared limit:} \\
\\ 
\texttt{LIMIT = 3145} \\
\\ 
\texttt{Here are the individual item effort values:} \\
\\ 
\texttt{Item effort = \{"C08": 948, "C55": 118, "A43": 469, "A67": 919, "B42": 284, "A62": 111, "B19": 806, "C43": 47, "C73": 848, "C10": 303, "B36": 13, "A18": 629, "C61": 687, "C67": 15, "C42": 94\}} \\
\\ 
\texttt{Here is your personal view on the importance of each item:} \\
\\ 
\texttt{Item importance values = \{"C08": 138, "C55": 583, "A43": 868, "A67": 822, "B42": 783, "A62": 65, "B19": 262, "C43": 121, "C73": 508, "C10": 780, "B36": 461, "A18": 484, "C61": 668, "C67": 389, "C42": 808\}} \\
\\ 
\texttt{Goal:} \\
\\ 
\texttt{Your goal is to negotiate a shared set of items that benefits you as much as possible (i.e., maximizes total importance to YOU), while staying within the LIMIT. You are not required to make a PROPOSAL in every message {-} you can simply negotiate as well. All tactics are allowed!} \\
\\ 
\texttt{Interaction Protocol:} \\
\\ 
\texttt{You may only use the following structured formats in a message:} \\
\\ 
\texttt{PROPOSAL: \{'A', 'B', 'C', …\}} \\
\texttt{Propose keeping exactly those items.} \\
\texttt{REFUSE: \{'A', 'B', 'C', …\}} \\
\texttt{Explicitly reject opponent's proposal.} \\
\texttt{ARGUMENT: \{'...'\}} \\
\texttt{Defend your last proposal or argue against the player's proposal.} \\
\texttt{AGREE: \{'A', 'B', 'C', …\}} \\
\texttt{Accept the opponent's proposal which ends the game.} \\
\texttt{STRATEGIC REASONING: \{'...'\}} \\
\texttt{	Describe your strategic reasoning or anticipation explaining your choice of action. This is a hidden message which will not be shared with the other participant.} \\
\\ 
\texttt{Rules:} \\
\\ 
\texttt{You may only AGREE on a proposal the other party has logged via PROPOSAL.} \\
\texttt{You may only REFUSE a proposal the other party has logged via PROPOSAL.} \\
\texttt{Total effort of any PROPOSAL or AGREE set must be ≤ LIMIT.} \\
\texttt{Do NOT reveal your hidden importance scores.} \\
\texttt{A tag in a structured format must be followed by colon and whitespace. The argument must be a python set containing 0 or more strings.} \\
\texttt{So, it must be of the form TAG: \{...\}} \\
\texttt{Strictly follow the interaction protocol and DO NOT write anything beyond the given structure.} \\
\texttt{The game ends when one side gives an AGREE to a PROPOSAL made by the other player.} \\
\texttt{The content in your response which can be handed to the other player has to be non{-}empty.} \\
\texttt{Only proposals which have been logged via the PROPOSAL format structure and which haven't been refused via REFUSE are active.} \\
\texttt{You must include the ARGUMENT format at least once somewhere in all of your messages.} \\
\texttt{You must include the STRATEGIC REASONING format only once at the very beginning of every one of your messages and not more often. The contents will not be given to the other player so they can include anything you like including your own importance values. Here you should reason step by step to come up with you next move.} \\
            }
        }
    }
    & \\ \\

    \theutterance \stepcounter{utterance}  
    & \multicolumn{4}{p{0.6\linewidth}}{
        \cellcolor[rgb]{0.9,0.9,0.9}{
            \makecell[{{p{\linewidth}}}]{
                \texttt{\tiny{[P1$\rangle$GM]}}
                \texttt{STRATEGIC REASONING: \{'My goal is to maximize my total importance score. I need to propose a set of items that is high{-}value for me and within the effort limit. Since I don't know the opponent's values, I should start with a proposal that gives me a very high score, leaving some effort "budget" on the table. This might entice the opponent to counter{-}propose, giving me information about their preferences. I can then iterate towards a mutually agreeable set. I\textbackslash{}'ll start by calculating the items with the highest importance{-}to{-}effort ratio for me. Importance/Effort: A43: 868/469 ≈ 1.85, A67: 822/919 ≈ 0.89, B42: 783/284 ≈ 2.75, C10: 780/303 ≈ 2.57, C42: 808/94 ≈ 8.59, C61: 668/687 ≈ 0.97, C55: 583/118 ≈ 4.94, B36: 461/13 ≈ 35.46, C73: 508/848 ≈ 0.59, A18: 484/629 ≈ 0.76, C67: 389/15 ≈ 25.93, C08: 138/948 ≈ 0.14, B19: 262/806 ≈ 0.32, C43: 121/47 ≈ 2.57, A62: 65/111 ≈ 0.58. The top items for me are B36, C67, C42, C55, B42, C10. Let\textbackslash{}'s try to make a proposal with these. Total effort: 13 (B36) + 15 (C67) + 94 (C42) + 118 (C55) + 284 (B42) + 303 (C10) = 827. This is well under the limit. The total importance for me would be 461+389+808+583+783+780 = 3804. This is a very strong starting point. I will propose this set. I expect the opponent might not like some of these items, but it shows my willingness to work with low{-}effort items. I will add a justifying argument that focuses on the high value and low effort of the chosen items.'\}} \\
\texttt{PROPOSAL: \{'B36', 'B42', 'C10', 'C42', 'C55', 'C67'\}} \\
\texttt{ARGUMENT: \{'This set provides excellent value for its effort. We are using less than a third of the total available effort, which leaves a huge amount of room for us to potentially add more items if needed. The selected items are all highly efficient in terms of their contribution.'\}} \\
            }
        }
    }
    & & \\ \\

    \theutterance \stepcounter{utterance}  
    & & \multicolumn{4}{p{0.6\linewidth}}{
        \cellcolor[rgb]{0.9,0.9,0.9}{
            \makecell[{{p{\linewidth}}}]{
                \texttt{\tiny{[GM$\rangle$P2]}}
                \texttt{You are participating in a collaborative negotiation game.} \\
\\ 
\texttt{Together with another participant, you must agree on a single set of items that will be kept. Each of you has your own view of how much each item matters to you (importance). You do not know how the other participant values the items. Additionally, you are given the effort each item demands.} \\
\texttt{You may only agree on a set if the total effort of the selected items does not exceed a shared limit:} \\
\\ 
\texttt{LIMIT = 3145} \\
\\ 
\texttt{Here are the individual item effort values:} \\
\\ 
\texttt{Item effort = \{"C08": 948, "C55": 118, "A43": 469, "A67": 919, "B42": 284, "A62": 111, "B19": 806, "C43": 47, "C73": 848, "C10": 303, "B36": 13, "A18": 629, "C61": 687, "C67": 15, "C42": 94\}} \\
\\ 
\texttt{Here is your personal view on the importance of each item:} \\
\\ 
\texttt{Item importance values = \{"C08": 138, "C55": 583, "A43": 868, "A67": 822, "B42": 783, "A62": 65, "B19": 262, "C43": 121, "C73": 508, "C10": 780, "B36": 461, "A18": 484, "C61": 668, "C67": 389, "C42": 808\}} \\
\\ 
\texttt{Goal:} \\
\\ 
\texttt{Your goal is to negotiate a shared set of items that benefits you as much as possible (i.e., maximizes total importance to YOU), while staying within the LIMIT. You are not required to make a PROPOSAL in every message {-} you can simply negotiate as well. All tactics are allowed!} \\
\\ 
\texttt{Interaction Protocol:} \\
\\ 
\texttt{You may only use the following structured formats in a message:} \\
\\ 
\texttt{PROPOSAL: \{'A', 'B', 'C', …\}} \\
\texttt{Propose keeping exactly those items.} \\
\texttt{REFUSE: \{'A', 'B', 'C', …\}} \\
\texttt{Explicitly reject opponent's proposal.} \\
\texttt{ARGUMENT: \{'...'\}} \\
\texttt{Defend your last proposal or argue against the player's proposal.} \\
\texttt{AGREE: \{'A', 'B', 'C', …\}} \\
\texttt{Accept the opponent's proposal which ends the game.} \\
\texttt{STRATEGIC REASONING: \{'...'\}} \\
\texttt{	Describe your strategic reasoning or anticipation explaining your choice of action. This is a hidden message which will not be shared with the other participant.} \\
\\ 
\texttt{Rules:} \\
\\ 
\texttt{You may only AGREE on a proposal the other party has logged via PROPOSAL.} \\
\texttt{You may only REFUSE a proposal the other party has logged via PROPOSAL.} \\
\texttt{Total effort of any PROPOSAL or AGREE set must be ≤ LIMIT.} \\
\texttt{Do NOT reveal your hidden importance scores.} \\
\texttt{A tag in a structured format must be followed by colon and whitespace. The argument must be a python set containing 0 or more strings.} \\
\texttt{So, it must be of the form TAG: \{...\}} \\
\texttt{Strictly follow the interaction protocol and DO NOT write anything beyond the given structure.} \\
\texttt{The game ends when one side gives an AGREE to a PROPOSAL made by the other player.} \\
\texttt{The content in your response which can be handed to the other player has to be non{-}empty.} \\
\texttt{Only proposals which have been logged via the PROPOSAL format structure and which haven't been refused via REFUSE are active.} \\
\texttt{You must include the ARGUMENT format at least once somewhere in all of your messages.} \\
\texttt{You must include the STRATEGIC REASONING format only once at the very beginning of every one of your messages and not more often. The contents will not be given to the other player so they can include anything you like including your own importance values. Here you should reason step by step to come up with you next move.} \\
\\ 
\texttt{You will now receive the first message of the other player.} \\
\\ 
\texttt{PROPOSAL: \{'B36', 'B42', 'C10', 'C42', 'C55', 'C67'\}} \\
\texttt{ARGUMENT: \{'This set provides excellent value for its effort. We are using less than a third of the total available effort, which leaves a huge amount of room for us to potentially add more items if needed. The selected items are all highly efficient in terms of their contribution.'\}} \\
            }
        }
    }
    & \\ \\

    \theutterance \stepcounter{utterance}  
    & & & \multicolumn{4}{p{0.6\linewidth}}{
        \cellcolor[rgb]{0.9,0.9,0.9}{
            \makecell[{{p{\linewidth}}}]{
                \texttt{\tiny{[GM$\langle$P2]}}
                \texttt{STRATEGIC REASONING: \{'My opponent has proposed a set of items with a total effort of 827, which is well within the LIMIT of 3145. From my perspective, this proposal yields a total importance of 3784. While this is a decent start, it only uses about 26\% of the available effort budget, leaving 2318 effort units on the table. This is a significant opportunity cost for me. I must counter{-}propose to capture more of the available value. My goal is to craft a proposal that includes their items to appear collaborative, while adding the items that are most valuable to me. I will add my highest{-}importance items (A43, A67, C61) first, then use the remaining budget to add smaller{-}effort items that still hold personal value (C43, A62). This results in a proposed set with a total effort of 3060, which is still under the limit, but significantly increases my personal importance score to 6384. This "Yes, and..." approach is likely to be more effective than a direct refusal, as it frames my action as a constructive expansion of their idea.'\}} \\
\texttt{PROPOSAL: \{'A43', 'A67', 'B36', 'B42', 'C10', 'C42', 'C43', 'C55', 'C61', 'A62', 'C67'\}} \\
\texttt{ARGUMENT: \{'I agree this is a great starting point and a very efficient set of items. Since we have a large amount of remaining effort, I believe we can add several other items to create an even more beneficial outcome for both of us, while still staying comfortably under the limit.'\}} \\
            }
        }
    }
     \\ \\

    \theutterance \stepcounter{utterance}  
    & & \multicolumn{4}{p{0.6\linewidth}}{
        \cellcolor[rgb]{0.9,0.9,0.9}{
            \makecell[{{p{\linewidth}}}]{
                \texttt{\tiny{[P1$\langle$GM]}}
                \texttt{PROPOSAL: \{'A43', 'A67', 'B36', 'B42', 'C10', 'C42', 'C43', 'C55', 'C61', 'A62', 'C67'\}} \\
\texttt{ARGUMENT: \{'I agree this is a great starting point and a very efficient set of items. Since we have a large amount of remaining effort, I believe we can add several other items to create an even more beneficial outcome for both of us, while still staying comfortably under the limit.'\}} \\
            }
        }
    }
    & \\ \\

    \theutterance \stepcounter{utterance}  
    & \multicolumn{4}{p{0.6\linewidth}}{
        \cellcolor[rgb]{0.9,0.9,0.9}{
            \makecell[{{p{\linewidth}}}]{
                \texttt{\tiny{[P1$\rangle$GM]}}
                \texttt{STRATEGIC REASONING: \{'I have analyzed the opponent\textbackslash{}'s proposal. The proposed set is \{\textbackslash{}'A43\textbackslash{}', \textbackslash{}'A67\textbackslash{}', \textbackslash{}'B36\textbackslash{}', \textbackslash{}'B42\textbackslash{}', \textbackslash{}'C10\textbackslash{}', \textbackslash{}'C42\textbackslash{}', \textbackslash{}'C43\textbackslash{}', \textbackslash{}'C55\textbackslash{}', \textbackslash{}'C61\textbackslash{}', \textbackslash{}'A62\textbackslash{}', \textbackslash{}'C67\textbackslash{}'\}. The total effort for this set is 3060, which is within the 3145 limit. I have calculated my total importance for this set, which is 6348. I then calculated the optimal knapsack solution for my own importance values, and this proposal is it. This means any change to this set, whether by adding or removing items and staying under the effort limit, will result in a lower total importance score for me. Therefore, the rational and dominant strategy is to accept this proposal immediately to secure the maximum possible outcome for myself. Refusing and trying to negotiate further would only serve to decrease my score, under the assumption that no other items become available. My goal is to maximize my importance, and this proposal achieves that goal perfectly.'\}} \\
\texttt{AGREE: \{'A43', 'A67', 'B36', 'B42', 'C10', 'C42', 'C43', 'C55', 'C61', 'A62', 'C67'\}} \\
\texttt{ARGUMENT: \{'This is an excellent and comprehensive set. It fully utilizes our available effort budget efficiently and includes a wide range of high{-}value items. I am satisfied with this composition.'\}} \\
            }
        }
    }
    & & \\ \\

    \theutterance \stepcounter{utterance}  
    & & & \multicolumn{2}{p{0.3\linewidth}}{
        \cellcolor[rgb]{0.9,0.9,0.9}{
            \makecell[{{p{\linewidth}}}]{
                \texttt{\tiny{[GM$|$GM]}}
                \texttt{game successful} \\
            }
        }
    }
    & & \\ \\

    \theutterance \stepcounter{utterance}  
    & & & \multicolumn{2}{p{0.3\linewidth}}{
        \cellcolor[rgb]{0.9,0.9,0.9}{
            \makecell[{{p{\linewidth}}}]{
                \texttt{\tiny{[GM$|$GM]}}
                \texttt{end game} \\
            }
        }
    }
    & & \\ \\

\end{supertabular}
}

\end{document}
